\chapter[Trabalhos Relacionados]{Trabalhos Relacionados}
\label{cap:relacionados}

Este capítulo situa o desenvolvimento do RECLin-PT em relação a recursos e estudos de
extração de informação clínica. Parte das campanhas em inglês que consolidaram a
extração de relações clínicas, passa pelos recursos e modelos para o reconhecimento
de entidades clínicas em português, incluindo antecedentes desenvolvidos no próprio
campus, e chega aos poucos trabalhos de extração de relações clínicas em português,
para identificar o que esses estudos oferecem ao modelo e o que deixam em aberto.

\section{Recursos e desafios em inglês}

A extração de relações clínicas consolidou-se em língua inglesa em torno de
campanhas de avaliação como as séries \foreignlanguage{english}{Informatics for
Integrating Biology and the Bedside} (i2b2) e \foreignlanguage{english}{National
NLP Clinical Challenges} (n2c2) \cite{uzuner_2011_i2b2}, que disponibilizaram
corpora anotados e tarefas compartilhadas para problemas como relações entre
medicamentos, dosagens e efeitos adversos. Esses esforços estabeleceram protocolos
de avaliação e evidenciaram a importância de tratar adequadamente o desbalanceamento
de classes e o vazamento de dados entre partições, preocupações transpostas para
este trabalho, respectivamente, pela geração controlada de pares candidatos
(Seção~\ref{sec:candidatos}) e pelo particionamento em nível de documento
(Seção~\ref{sec:particao}).

Quanto aos modelos, o BioBERT \cite{lee_2020_biobert} demonstrou os
ganhos da adaptação de domínio biomédico em inglês, estabelecendo o precedente que
torna a escolha entre \emph{encoders} de domínio e gerais uma decisão a examinar no
desenvolvimento do RECLin-PT. Quanto à representação de entrada,
a delimitação do par por marcadores de entidade inseridos no próprio texto
\cite{soares_2019_markers} tornou-se a técnica padrão para a extração de relações com modelos baseados em \foreignlanguage{english}{BERT}, sendo também adotada neste trabalho
(Seção~\ref{sec:markers}).

\section{Reconhecimento de entidades clínicas em português}
\label{sec:recursos_pt}

O desenvolvimento de PLN clínico em português foi por muito tempo limitado pela
ausência de corpora anotados de acesso controlado. O \textbf{SemClinBr}
\cite{oliveira_2022_semclinbr} preencheu parte dessa lacuna ao disponibilizar mil
documentos clínicos de múltiplas instituições e especialidades, anotados
semanticamente e acompanhados de dicionários de abreviações e de pistas de negação.
Para este trabalho, sua contribuição central é a anotação de relações, inclusive de
negação, entre entidades. Como o corpus anota apenas relações positivas, usá-lo para
classificar todos os pares de um documento exige construir os exemplos negativos
(Seção~\ref{sec:candidatos}). É o corpus sobre o qual o RECLin-PT é desenvolvido.

Os \emph{encoders} disponíveis para o português foram avaliados, até aqui, sobretudo
em reconhecimento de entidades. O BioBERTpt \cite{schneider_2020_biobertpt}
adaptou modelos BERT a narrativas clínicas e textos biomédicos em português
brasileiro e relatou ganhos em NER em dois corpora de narrativas clínicas. O
BERTimbau \cite{souza_2020_bertimbau}, referência de domínio geral, foi avaliado em
similaridade textual, inferência textual e reconhecimento de entidades, todas tarefas
de domínio geral. Nenhum dos dois trabalhos mede o efeito do pré-treinamento sobre a
extração de relações. O estudo mais próximo, de \citeonline{almeida_2026_benchmarkner},
compara sobre o SemClinBr, com hiperparâmetros comuns sempre que possível, variantes
do BioBERTpt, o BERTimbau, modelos BERT recentes e modelos de linguagem de grande
porte, também em reconhecimento de entidades. Nesse estudo, o melhor desempenho
coube a um modelo multilíngue de domínio geral (mmBERT-base), à frente das variantes
do BioBERTpt, o que indica que o domínio do pré-treino não é o único fator relevante
mesmo em NER. Esses resultados motivam a comparação entre \emph{encoders} antes da escolha daquele que será utilizado no RECLin-PT, mas não dispensam essa análise. Permanece em aberto a questão de saber se a vantagem observada no reconhecimento de entidades também se manifesta na extração de relações.

No reconhecimento de entidades clínicas em português, destacam-se também pesquisas
desenvolvidas no Instituto Federal do Espírito Santo, Campus Cachoeiro de Itapemirim.
\citeonline{torres_2024_ner} desenvolveu um modelo baseado na biblioteca spaCy para
identificar entidades clínicas em relatos de casos, inicialmente com relatos anotados
manualmente e, em seguida, com amostras do SemClinBr, para as quais definiu categorias
próprias de representação das informações médicas. Nos relatos não usados no treino,
os melhores resultados foram os das categorias Conclusão e Procedimento Diagnóstico,
com F1 de $0{,}71$ e $0{,}65$. Posteriormente, \citeonline{ramos_2025_clinpthyb}
propôs o ClinptHyb, abordagem híbrida que combina um modelo NER treinado com o spaCy
sobre o SemClinBr e um modelo de linguagem de grande porte
(\foreignlanguage{english}{Large Language Model}, LLM), utilizado como revisor das
anotações produzidas pelo primeiro. Com o ClinptHyb avaliado em uma amostra
estratificada de 33 relatos do conjunto de validação, o trabalho relata F1 superior ao
do modelo spaCy isolado em todas as categorias consideradas, incluindo $0{,}972$ na
categoria \texttt{Negation}. Ao
revisar as anotações, o trabalho também identificou variações na delimitação de
entidades e omissões de anotação no próprio SemClinBr.

Embora adotem estratégias distintas, os dois trabalhos concentram-se na
identificação e classificação de entidades clínicas. O F1 de $0{,}972$ em
\texttt{Negation}, por exemplo, mede o reconhecimento das expressões de negação, e não
a ligação entre cada expressão e o conceito negado, que é o objeto da relação
\texttt{negation\_of} (Seção~\ref{sec:re_tarefa}). As duas pesquisas apontam a
extração de relações entre as entidades reconhecidas como possibilidade de
continuidade, sem realizá-la. O RECLin-PT dá um passo nessa direção ao investigar
diretamente a extração de relações entre entidades clínicas, a partir das entidades
já anotadas no SemClinBr e com ênfase nas relações de negação. A integração entre um
reconhecimento automático de entidades e a extração de relações permanece fora do
escopo deste trabalho.

\section{Extração de relações clínicas em português}

A extração de relações clínicas em português é pouco representada na literatura
revisada. Uma revisão sistemática de 28 estudos de modelos de linguagem em português
voltados à saúde, publicados entre janeiro de 2020 e novembro de 2024, inventaria
classificação (16 modelos), reconhecimento de entidades nomeadas (8), modelagem de
tópicos, resposta a perguntas, geração de texto e agentes conversacionais, sem
registrar extração de relações \cite{shimaoka_2026_revisao}. O desequilíbrio não é
exclusivo do português: uma revisão de 94 estudos de extração de informação em texto
clínico registra 63 dedicados a reconhecimento de entidades contra 13 dedicados a
extração de relações \cite{navarro_2023_revisao}. O trabalho de extração de relações
em texto clínico em português aqui localizado \cite{silva_2023_oncologia} contribui
com um corpus próprio de oncologia em português brasileiro e menciona três
\emph{encoders} candidatos (BERT multilíngue, BioBERTpt e BERTimbau), mas não
reporta comparação entre eles e resolve a etapa de relações com uma rede BiLSTM, sem
teste de significância. O efeito do pré-treinamento de domínio sobre a extração de
relações, portanto, não é isolado ali.

\section{Síntese e posicionamento}

Os trabalhos revisados oferecem o corpus, os \emph{encoders} e as técnicas de
representação do par em que o RECLin-PT se apoia, e os antecedentes do campus
mostram uma linha de pesquisa em reconhecimento de entidades clínicas em português
que já relata resultados altos para as expressões de negação. Permanece, porém, a lacuna referente à etapa seguinte, que consiste em investigar a extração de relações entre entidades clínicas em português, especialmente as relações de negação, responsáveis por identificar os conceitos aos quais a negação está associada. Não foram identificados estudos que verifiquem, para o português, se a vantagem do
pré-treinamento clínico relatada em NER se transfere para a extração de relações, nem
que compare, sobre um mesmo corpus clínico em português, estratégias de modelagem
distintas sob protocolo único.

O RECLin-PT se posiciona nessa lacuna. É desenvolvido sobre o SemClinBr com foco na
negação clínica, e sua arquitetura e as estratégias investigadas combinam técnicas já
descritas na literatura, como os marcadores de entidade e a decomposição da decisão.
Sua contribuição central é o próprio modelo. As comparações de \emph{encoders} e de
estratégias, conduzidas sob as mesmas condições e com uma cadeia reprodutível, servem
para orientar sua configuração e, ao mesmo tempo, produzem evidência sobre a lacuna
descrita, que pode servir de referência para trabalhos futuros sobre o mesmo corpus.
